\chapter{FLC Identification Framework}
\label{chap:identification_framework}

\section{Motivation and Scope}

The identification of forming-limit points is a key step in the construction
of forming limit curves. After a forming test, the measured strain field must
be evaluated in order to determine the strain state at the onset of localized
necking. Since different identification criteria can lead to different limit
strains, it is useful to compare them within the same post-processing workflow.

In this project, the starting point was an existing graphical user interface
developed in the laboratory for the post-processing of DIC measurements. This
tool already included data import, visualization, plotting routines, and the
DIN/ISO~12004-2 and Volk–Hora identification methods. The contribution of the
present work was therefore to study an additional criterion based on the work
of Min, Stoughton, Carsley, and Lin, and to prepare a draft methodology for its
future integration into the existing framework.

\section{Existing Identification Criteria}

The existing framework contains two forming-limit identification methods. The
DIN/ISO~12004-2 method evaluates the strain distribution along a section
crossing the localization or crack region and determines the limit strain from
a spatial fit. The Volk–Hora criterion uses the temporal evolution of the
strain field to detect the onset of instability.
Both methods are available within the same GUI environment, which allows them
to be applied to the same DIC dataset.

\section{Draft of the Min–Stoughton Criterion}
\label{sec:minstoughton_draft}

The Min–Stoughton criterion was studied as a possible extension of the
existing FLC identification framework. The physical principle of the criterion,
namely the detection of localized necking from the evolution of the surface
curvature, is described in Section~\ref{sec:minstoughton_method}. The present
section concerns its translation into a working post-processing procedure for
the laboratory DIC data.

For Marciniak tests, this curvature is fitted on a small narrow surface patch
around the neck band, rather than on one manually selected line. This makes the
criterion less dependent on the exact line chosen by the operator. For
Nakazima tests, the same idea is applied after a geometric transformation,
because the sheet wraps around a hemispherical punch. The measured surface is
expressed relative to the punch geometry, so that the local curvature increase
caused by necking can be separated from the global curvature imposed by the
punch.

Within this project the Min–Stoughton method was therefore treated as a
methodology and data-integration task rather than as a finalized GUI feature.
The required VIC3D quantities are the measured surface positions and
displacements $(X,Y,Z,U,V,W)$, the strain fields $(\varepsilon_1,
\varepsilon_2)$, the crack or localization location used to place the point
matrix, and metadata such as punch diameter and sheet thickness. A provisional
standalone workflow was created to test these data links and to explore the
Nakazima transformation, but the detected onset remained sensitive to the
matrix dimensions, reference-frame choice, smoothing, and persistence
definition. The final integration therefore still requires a validated
parameter choice and a consistent treatment of decimated VIC3D frame sequences.

\subsection{Standalone Nakazima Workflow Used for Assessment}
\label{sec:minstoughton_workflow_assessment}

The standalone Nakazima implementation was used as an assessment workflow for
the present report. It does not replace the existing DIN/ISO and Volk--Hora GUI
tools, but it connects the mathematical method of
Section~\ref{sec:minstoughton_method} to the available laboratory DIC exports.
The workflow imports the VIC3D reference coordinates, displacement fields, and
principal strains, reconstructs the deformed surface coordinates, and reads the
fixed specimen metadata needed for the Nakazima transformation.

The user-drawn crack line is used to orient the point matrix. Its centre $O$ is
then selected from the local strain field as the maximum of the smoothed major
strain in a band around that line, so the matrix is not blindly centred on the
crack-line midpoint. The requested physical matrix dimensions $W_X$ and $W_Y$
are converted into the point counts $N_X$ and $N_Y$ using the measured DIC grid
spacing according to Eq.~\eqref{eq:ms_matrix_counts}. The same DIC point IDs are
then followed through the recorded frames.

For the Nakazima branch, the implementation follows the coordinate definitions
of Eqs.~\eqref{eq:ms_punch_centre}--\eqref{eq:ms_rout}. The coordinate $D$ is
computed as the cumulative three-dimensional distance along the ordered
deformed DIC points of each matrix column, and $R_\mathrm{out}$ is computed from
the measured deformed point to the reconstructed punch centre. No ideal
spherical projection is applied, so the measured dimple remains present in the
radial perturbation. The reference subtraction, SAC addition, circle fitting,
curvature prediction, and onset criterion then follow
Eqs.~\eqref{eq:ms_reference_subtraction}--\eqref{eq:ms_delta}.

The reference frame $M$ is handled as a time-based choice rather than a silent
frame-count fraction, because the VIC3D sequences used here have different
frame rates before and near fracture. The selected reference frame ID and time
are stored with the output so that the subtraction state is traceable.

The implementation therefore relies on a small set of constants that originate
from four distinct sources. A first group is fixed by the method of Min et
al.~\cite{Min2017}. A second group is read from the specimen metadata. A third
group is derived from the measured DIC grid for each specimen. The remaining
constants are the analysis parameters chosen by the user and examined in the
sensitivity study. These constants and their origin are summarised in
Table~\ref{tab:ms_constants}.

\begin{table}[H]
    \centering
    \small
    \caption{Constants required by the Min--Stoughton workflow and their origin.}
    \label{tab:ms_constants}
    \begin{tabular}{l p{0.40\linewidth} p{0.30\linewidth}}
        \hline
        Symbol & Meaning & Origin \\
        \hline
        \multicolumn{3}{l}{\textit{Fixed by the method}}\\
        $n$ & Number of consecutive frames required for onset & Min et al.~\cite{Min2017}, $n=8$ \\
        $\alpha$ & Threshold multiplier, $\Delta=\alpha\,\mathrm{SAC}$ & Paper default $\alpha=1/10$ \\
        \hline
        \multicolumn{3}{l}{\textit{Read from specimen metadata}}\\
        $r_N$ & Punch radius, used to reconstruct the punch centre & \texttt{sample\_ID.xml} / VIC3D project \\
        $t_0$ & Initial sheet thickness & \texttt{sample\_ID.xml} / VIC3D project \\
        \hline
        \multicolumn{3}{l}{\textit{Derived from the DIC measurement}}\\
        $d_X,\,d_Y$ & DIC grid spacing across and along the band & Median spacing of adjacent DIC points \\
        $N_X,\,N_Y$ & Matrix point counts & Computed from Eq.~\eqref{eq:ms_matrix_counts} \\
        $O$ & Neck centre that anchors the matrix & Local maximum of smoothed $\varepsilon_1$ near the crack line \\
        \hline
        \multicolumn{3}{l}{\textit{User-selected analysis parameters}}\\
        $W_X,\,W_Y$ & Physical matrix width and length & Chosen, studied in Section~\ref{sec:results_task2} \\
        $\mathrm{SAC}$ & Superimposed artificial curvature & Chosen value, e.g. $5\times10^{-4}\,\mathrm{mm}^{-1}$ \\
        $M$ & Reference frame for background subtraction & Time-based choice, e.g. $M=0.8$ \\
        \hline
    \end{tabular}
\end{table}

This section defines how the criterion is assembled. The numerical sensitivity
of the current implementation, and its comparison with the existing DIN/ISO and
Volk--Hora outputs, are reported separately in Section~\ref{sec:results_task2}.
